
\section{计数\BookRecordIndex{T0335}伪全纯图像}\label{sec:7-3}
\zhindex{jishuweiquanchuntuxiang}{计数伪全纯图像}

本节考察紧\BookRecordIndex{T0058}半正\BookRecordIndex{T0358}辛流形 $(M,\omega)$ 上与\BookRecordIndex{T0088}殆复结构光滑族
\[
  S^2\longrightarrow\mathcal J_\tau(M,\omega),
  \qquad z\longmapsto J_z
\]
相关联的\BookRecordIndex{T0143}格罗莫夫--威滕不变量，正如第~\ref{sec:6-7}节所述。固定常数
$\kappa>0$，并回忆
\[
  \mathcal J_+(M,\omega;\kappa)\subset\mathcal J_\tau(M,\omega)
\]
表示 $M$ 上所有满足如下条件的被 $\omega$ \BookRecordIndex{T0372}驯服的\BookRecordIndex{T0088}殆复结构 $J$ 所成的空间：每个能量
\[
  E(u)=\int_{S^2}u^*\omega\le\kappa
\]
的 $J$-全纯球面 $u:S^2\to M$ 都有非负 Chern 数。由引理~\ref{lem:6-4-8}，这是
$\mathcal J_\tau(M,\omega)$ 的一个非空、道路连通的开子集。还要回忆，
\[
  \mathcal J_+(S^2;M,\omega;\kappa)
  \subset C^\infty\bigl(S^2,\mathcal J_+(M,\omega;\kappa)\bigr)
\]
表示所有可缩光滑映射
\[
  S^2\longrightarrow\mathcal J_+(M,\omega;\kappa),
  \qquad z\longmapsto J_z
\]
所成的集合。这同样是
$C^\infty(S^2,\mathcal J_\tau(M,\omega))$ 的一个非空、道路连通的开子集。在几何上，
可以把 $\mathcal J_+(S^2;M,\omega;\kappa)$ 看成所有充分接近常值族的\BookRecordIndex{T0088}殆复结构光滑族所成的
集合。若要改用 $M$ 上任意的被 $\omega$ \BookRecordIndex{T0372}驯服的\BookRecordIndex{T0088}殆复结构族，就必须施加一个比\BookRecordIndex{T0058}半正性更强的
条件，即对每个 $A\in H_2(M;\mathbb Z)$，
\[
  \omega(A)>0,\qquad c_1(A)\ge 2-n
  \qquad\Longrightarrow\qquad c_1(A)\ge0.
\]

现在固定非负整数 $k$、指标集 $I\subset\{1,\ldots,k\}$，以及 $S^2$ 上两两不同的点所成的
元组 $\mathbf w=\{w_i\}_{i\in I}$。令
\[
  \mathbf J=\{J_z\}_{z\in S^2}
  \in\mathcal J_{\mathrm{reg}}(S^2;M,\omega;\mathbf w,\kappa)
  :=\mathcal J_{\mathrm{reg}}(S^2;M,\omega;\mathbf w)
    \cap\mathcal J_+(S^2;M,\omega;\kappa)
\]
为第~\ref{sec:6-7}节（见定义~\ref{def:6-7-10}）所定义的\BookRecordIndex{T0088}殆复结构正则族。设
$A\in H_2(M;\mathbb Z)$ 满足 $\omega(A)\le\kappa$。于是，类 $A$ 中的
$\{J_z\}$-全纯球面 $u:S^2\to M$ 所成的\BookRecordIndex{T0246}模空间
\[
  \mathcal M(A;\mathbf J)
  :=\bigl\{u:S^2\to M\mid
  du(z)\circ j=J_z(u(z))\circ du(z),\ [u]=A\bigr\}
\]
是维数 $2n+2c_1(A)$ 的光滑流形。相应的带标记球面\BookRecordIndex{T0246}模空间
\[
  \mathcal M_{0,k}(A;\mathbf w,\mathbf J)
  :=\bigl\{(u,\mathbf z)\in\mathcal M(A;\mathbf J)\times(S^2)^k
  \mid z_i\ne z_j\ (i\ne j),\ z_i=w_i\ \forall i\in I\bigr\}
\]
的维数是
\[
  \mu(A,k,I):=2n+2c_1(A)+2(k-\#I).
\]
定理~\ref{thm:6-7-1} 说明，只要
$\mathbf J\in\mathcal J_{\mathrm{reg}}(S^2;M,\omega;\mathbf w,\kappa)$ 且
$\omega(A)\le\kappa$，\BookRecordIndex{T0283}求值映射
\[
  \operatorname{ev}_{\mathbf w,\mathbf J}:
  \mathcal M_{0,k}(A;\mathbf w,\mathbf J)\longrightarrow M^k
\]
就是\BookRecordIndex{T0337}伪循环。这便如下给出了\BookRecordIndex{T0335}伪全纯图像的\BookRecordIndex{T0143}格罗莫夫--威滕不变量。记
$\pi_i:M^k\to M$ 为到第 $i$ 个因子的投影。

\begin{theorem}\label{thm:7-3-1}
设 $(M,\omega)$ 是闭\BookRecordIndex{T0058}半正辛 $2n$-流形，$\kappa>0$，且球面同调类
$A\in H_2(M;\mathbb Z)$ 满足 $\omega(A)\le\kappa$。设 $k$ 是非负整数，
$I\subset\{1,\ldots,k\}$ 是指标集，$\mathbf w=\{w_i\}_{i\in I}$ 是 $S^2$ 上两两不同的
点所成的元组。设
$\mathbf J\in\mathcal J_{\mathrm{reg}}(S^2;M,\omega;\mathbf w,\kappa)$。则对
$a_i\in H^*(M)$，取对偶于
$\pi_1^*a_1\smile\cdots\smile\pi_k^*a_k$ 的\BookRecordIndex{T0337}伪循环 $f:U\to M^k$，由
\begin{equation}\label{eq:7-3-1}
  \widetilde{\mathrm{GW}}^{M,I}_{A,k}(a_1,\ldots,a_k)
  :=f\mathbin{\boldsymbol\cdot}\operatorname{ev}_{\mathbf w,\mathbf J}
\end{equation}
定义的同态
\[
  \widetilde{\mathrm{GW}}^{M,I}_{A,k}:
  H^*(M)^{\otimes k}\longrightarrow\mathbb Z
\]
与\BookRecordIndex{T0150}固定标记点元组 $\mathbf w$、正则\BookRecordIndex{T0088}殆复结构族
$\mathbf J\in\mathcal J_{\mathrm{reg}}(S^2;M,\omega;\mathbf w,\kappa)$ 以及定义中所用的
\BookRecordIndex{T0337}伪循环 $f$ 都无关。
\end{theorem}

\begin{proof}
由定理~\ref{thm:6-7-1}，$\operatorname{ev}_{\mathbf w,\mathbf J}$ 的配边类与
$\mathbf w$ 和 $\mathbf J$ 无关。因此，由引理~\ref{lem:6-5-5}，交数
$f\mathbin{\boldsymbol\cdot}\operatorname{ev}_{\mathbf w,\mathbf J}$ 与
$\mathbf w$ 和 $\mathbf J$ 无关。再由引理~\ref{lem:6-5-7}，它只依赖于\BookRecordIndex{T0337}伪循环 $f$ 所代表的
上同调类。这就证明了定理。
\end{proof}

式~\eqref{eq:7-3-1} 定义的不变量称为\emph{类 $A$ 中、由 $I$ 标号的\BookRecordIndex{T0066}标记点固定的
$k$-点图像\BookRecordIndex{T0143}格罗莫夫--威滕不变量}。它在几何上可作如下解释。选取一般位置的循环 $X_i$，
使其 Poincaré 对偶为 $a_i$；再固定 $S^2$ 上两两不同的点所成的元组
$\mathbf w=\{w_i\}_{i\in I}$。那么，式~\eqref{eq:7-3-1} 带符号计数元组
\[
  (u,\mathbf z)=\bigl(u,\{z_i\}_{i\notin I}\bigr),
\]
其中 $u:S^2\to M$ 是代表类 $A$ 的 $\{J_z\}$-全纯球面，而对
$i\in\{1,\ldots,k\}\setminus I$，各点 $z_i$ 在
$S^2\setminus\{w_i\mid i\in I\}$ 中两两不同；所有\BookRecordIndex{T0066}标记点都被映到相应的循环，即
\[
  u(w_i)\in X_i\quad(i\in I),
  \qquad
  u(z_i)\in X_i\quad(i\notin I).
\]
这确实可良定义为\BookRecordIndex{T0337}伪循环的交数，恰是定理~\ref{thm:7-3-1} 的断言。我们强调这一定义的如下
特点。
\begin{itemize}
\item 定义该不变量时，可以假设\BookRecordIndex{T0337}伪循环 $f$ 实际上是定向子流形
$X=X_1\times\cdots\times X_k\subset M^k$。代价是必须把上同调类 $a_i$ 换成适当的整数倍；
若要得到整数值不变量，则可能需要使用 $a_i$ 的更一般代表。
\item 不变量 $\widetilde{\mathrm{GW}}^{M,I}_{A,k}(a_1,\ldots,a_k)$ 只有在各类 $a_i$
满足维数条件
\begin{equation}\label{eq:7-3-2}
  \sum_{i=1}^k\deg(a_i)=2n+2c_1(A)+2k-2\#I
\end{equation}
时才可能非零。
\item 当 $\#I<3$ 时，$\widetilde{\mathrm{GW}}^{M,I}_{A,k}$ 恒为零，除非
$A=0$ 且 $k=\#I$（见习题~\ref{ex:7-3-4}）。这是扰动会破坏对称性的结果。若用
$G=\operatorname{PSL}(2,\mathbb C)$ 的元素重参数化一个 $\{J_z\}$-全纯球面，它便不再满足
同一个方程。因此，我们不能对 $G$ 的作用取商，以把\BookRecordIndex{T0246}模空间降到其“正确”维数。这一点也
体现在我们的设置中：除非 $u$ 为常值（从而 $A=0$），否则 $u$ 与 $u\circ\phi$ 的图像彼此
不同。
\item 若 $\#I=3$ 且 $A$ 满足式~\eqref{eq:6-6-1}，则不变量
$\widetilde{\mathrm{GW}}^{M,I}_{A,k}$ 与 $\mathrm{GW}^M_{A,k}$ 相同（见
命题~\ref{prop:7-3-2}）。
\end{itemize}

\begin{proposition}\label{prop:7-3-2}
设 $(M,\omega)$ 是紧\BookRecordIndex{T0058}半正\BookRecordIndex{T0358}辛流形，$k\ge3$ 是整数，并假设
$A\in H_2(M;\mathbb Z)$ 满足式~\eqref{eq:6-6-1}。则
\[
  \widetilde{\mathrm{GW}}^{M,\{1,2,3\}}_{A,k}
  =\mathrm{GW}^M_{A,k}:H^*(M)^{\otimes k}\longrightarrow\mathbb Z.
\]
\end{proposition}

\begin{proof}
这直接来自定理~\ref{thm:6-7-1}\textup{(iii)}。
\end{proof}

\begin{exercise}\label{ex:7-3-3}
设 $(M,\omega)$ 单调，最小 Chern 数为 $N$。设 $k$ 是正整数，并设指标集
$I\subset\{1,\ldots,k\}$ 满足 $3\le\#I\le N+2$。固定同调类
$A\in H_2(M;\mathbb Z)$，并设 $a_1,\ldots,a_k\in H^*(M)$ 是一族满足
式~\eqref{eq:7-3-2} 的上同调类。证明
\[
  \widetilde{\mathrm{GW}}^{M,I}_{A,k}(a_1,\ldots,a_k)
  =\mathrm{GW}^{M,I}_{A,k}(a_1,\ldots,a_k).
\]
\end{exercise}

\begin{exercise}\label{ex:7-3-4}
假设 $\#I<3$。证明 $\widetilde{\mathrm{GW}}^{M,I}_{A,k}=0$，除非 $A=0$ 且
$k=\#I$。还要证明
\[
  \widetilde{\mathrm{GW}}^{M,\{1\}}_{0,1}(a)=\int_M a,
  \qquad
  \widetilde{\mathrm{GW}}^{M,\{1,2\}}_{0,2}(a_1,a_2)
  =\int_M a_1\smile a_2.
\]

\emph{提示：}为证明第一个断言，选取一列收敛到与 $z$ 无关的\BookRecordIndex{T0088}殆复结构
$J\in\mathcal J_{\mathrm{reg}}(M,\omega)$ 的\BookRecordIndex{T0088}殆复结构族 $\{J_z^\nu\}$。利用维数说明，不存在
代表类 $A$、并穿过与 $a_i$ 对偶的一般循环的 $J$-\BookRecordIndex{T0340}稳定映射。为证明第二个断言，利用能量
说明式~\eqref{eq:6-7-1} 在类 $A=0$ 中的解只能是常值映射。由
引理~\ref{lem:6-7-6}，常值解是正则的。
\end{exercise}

\begin{example}[杯积]\label{ex:7-3-5}
当 $A=0$ 时，不变量 $\widetilde{\mathrm{GW}}$ 由杯积给出。更确切地说，只要
$a_i\in H^*(M)$ 满足 $\sum_i\deg(a_i)=2n$，就有
\[
  \widetilde{\mathrm{GW}}^{M,\{1,\ldots,k\}}_{0,k}(a_1,\ldots,a_k)
  =\int_M a_1\smile\cdots\smile a_k.
\]
事实上，每个代表类 $A=0$ 的 $\{J_z\}$-全纯球面都是常值的，因而\BookRecordIndex{T0246}模空间
$\mathcal M(0,\mathbf w;\mathbf J)$ 微分同胚于 $M$。由习题~\ref{ex:7-3-4}，这个公式对
所有 $k\ge1$ 都成立。另请比较例~\ref{ex:7-1-3}。
\end{example}

下一个命题把定理~\ref{thm:7-3-1} 的不变量 $\widetilde{\mathrm{GW}}$ 与乘积流形
$\widetilde M:=S^2\times M$ 的\BookRecordIndex{T0143}格罗莫夫--威滕不变量作比较，其中假设
$\widetilde M$ 本身\BookRecordIndex{T0058}半正。

\begin{proposition}\label{prop:7-3-6}
假设乘积流形 $\widetilde M:=S^2\times M$ \BookRecordIndex{T0058}半正。设
$A\in H_2(M;\mathbb Z)$ 是球面同调类，并记
\[
  \widetilde A:=[S^2\times\mathrm{pt}]+\iota_*A
  \in H_2(\widetilde M;\mathbb Z),
\]
其中 $\iota:M\to S^2\times M$ 是纤维的包含映射。则对 $a_i\in H^*(M)$，有
\[
  \widetilde{\mathrm{GW}}^{M,I}_{A,k}(a_1,\ldots,a_k)
  =\mathrm{GW}^{\widetilde M}_{\widetilde A,k}
  (\widetilde a_1,\ldots,\widetilde a_k),
\]
其中
\[
  \widetilde a_i:=
  \begin{cases}
    \operatorname{pr}^*a_i\smile\pi^*\sigma,&i\in I,\\
    \operatorname{pr}^*a_i,&i\notin I.
  \end{cases}
\]
这里 $\operatorname{pr}:\widetilde M\to M$ 是到第二个因子的投影，
$\pi:\widetilde M\to S^2$ 是到第一个因子的投影，而 $\sigma$ 是 $H^2(S^2)$ 的正生成元。
\end{proposition}

\begin{proof}
假设各类 $a_i$ 满足式~\eqref{eq:7-3-2}。则各类 $\widetilde a_i$ 满足把 $M,A$ 分别换成
$\widetilde M,\widetilde A$ 后的式~\eqref{eq:7-1-2}：
\[
  \sum_{i=1}^k\deg(\widetilde a_i)
  =\dim M+2c_1(A)+2k
  =\dim\widetilde M+2c_1(\widetilde A)+2k-6.
\]
条件 $\pi_*\widetilde A=[S^2]$ 保证 $\widetilde A$ 不是任何其他同调类的非平凡整数倍，
因而满足式~\eqref{eq:6-6-1}，可以按定理~\ref{thm:7-1-1} 定义亏格零\BookRecordIndex{T0143}格罗莫夫--威滕
不变量。

固定 $S^2$ 上两两不同的点所成的元组 $\mathbf w=\{w_i\}_{i\in I}$，并取常数
$\kappa\ge\omega(A)$。令
\[
  \mathbf J=\{J_z\}_{z\in S^2}
  \in\mathcal J_{\mathrm{reg}}(S^2;M,\omega;\mathbf w,\kappa).
\]
由 $\mathbf J$ 和 $S^2$ 上的标准复结构 $j$ 决定的乘积\BookRecordIndex{T0088}殆复结构
$\widetilde J\in\mathcal J(\widetilde M)$ 被\BookRecordIndex{T0363}辛形式
\[
  \widetilde\omega
  :=\operatorname{pr}^*\omega+\pi^*\operatorname{dvol}_{S^2}
  \in\Omega^2(\widetilde M)
\]
\BookRecordIndex{T0372}驯服。\sourcefootnote{原书称 $\widetilde J$ 与 $\widetilde\omega$ \BookRecordIndex{T0348}相容；但本段只假设
每个 $J_z$ 被 $\omega$ \BookRecordIndex{T0372}驯服，这只能\BookRecordIndex{T0326}推出乘积结构被 $\widetilde\omega$ \BookRecordIndex{T0372}驯服。\BookRecordIndex{T0348}相容性还需
额外假设每个 $J_z$ 与 $\omega$ \BookRecordIndex{T0348}相容。}此外，由正则族
$\mathbf J=\{J_z\}_{z\in S^2}$ 的定义可知：对每棵带\BookRecordIndex{T0321}特殊顶点
$0$ 的树 $T$，以及每族满足
\[
  \pi_*\widetilde A_0=[S^2],
  \qquad
  \pi_*\widetilde A_\alpha=0\quad(\alpha\in T\setminus\{0\})
\]
的同调类 $\{\widetilde A_\alpha\}_{\alpha\in T}\subset H_2(\widetilde M;\mathbb Z)$，
$\widetilde J$ 都满足定义~\ref{def:6-2-1} 的要求。尽管这不一定意味着
$\widetilde J\in\mathcal J_{\mathrm{reg}}(\widetilde M,\widetilde\omega)$，但已经足以用
$\widetilde J$ 构造定理~\ref{thm:7-1-1} 中类 $\widetilde A$ 的不变量。

现在假设对 $i=1,\ldots,k$，$a_i$ 是定向子流形 $X_i\subset M$ 的 Poincaré 对偶。则
$\widetilde a_i$ 对偶于子流形
\[
  \widetilde X_i:=
  \begin{cases}
    \{w_i\}\times X_i,&i\in I,\\
    S^2\times X_i,&i\notin I.
  \end{cases}
\]
此外，每条代表同调类 $\widetilde A_0$ 且满足
$\pi_*\widetilde A_0=[S^2]$ 的 $\widetilde J$-全纯曲线
$\widetilde u_0:S^2\to\widetilde M$，都能唯一地重参数化成 $\widetilde M$ 的
$\widetilde J$-全纯截面。因此，
$\mathrm{GW}^{\widetilde M}_{\widetilde A,k}(\widetilde a_1,\ldots,\widetilde a_k)$
计数元组 $\bigl(u,\{z_i\}_{i\notin I}\bigr)$，其中 $u:S^2\to M$ 是代表类 $A$ 的
$\{J_z\}$-全纯球面，且
\[
  u(w_i)\in X_i\quad(i\in I),
  \qquad
  u(z_i)\in X_i\quad(i\notin I).
\]
两个不变量相等，是因为具有这些性质的每个元组
$\bigl(u,\{z_i\}_{i\notin I}\bigr)$ 对两个不变量的贡献符号相同。这就证明了命题。
\end{proof}

命题~\ref{prop:7-3-2} 和习题~\ref{ex:7-3-3} 提示了下面的定义：对\BookRecordIndex{T0058}半正\BookRecordIndex{T0358}辛流形
$(M,\omega)$ 的任意同调类 $A\in H_2(M;\mathbb Z)$，无论它是否满足
式~\eqref{eq:6-6-1}，都可定义其\BookRecordIndex{T0143}格罗莫夫--威滕不变量。此外，这一定义允许任取\BookRecordIndex{T0150}固定标记点
集 $I\subset\{1,\ldots,k\}$，从而去除了式~\eqref{eq:7-2-10} 中
$\mathrm{GW}^{M,I}_{A,k}$ 定义里的那些很强的限制。

\begin{definition}\label{def:7-3-7}
设 $(M,\omega)$ 是闭\BookRecordIndex{T0058}半正\BookRecordIndex{T0358}辛流形，$k\ge3$ 是整数，
$I\subset\{1,\ldots,k\}$ 满足 $\#I\ge3$，且 $A\in H_2(M;\mathbb Z)$。定义 $M$ 中类
$A$ 的带 $k$ 个\BookRecordIndex{T0066}标记点的\emph{亏格零\BookRecordIndex{T0143}格罗莫夫--威滕不变量}为同态
\[
  \mathrm{GW}^M_{A,k}
  :=\widetilde{\mathrm{GW}}^{M,\{1,2,3\}}_{A,k}:
  H^*(M)^{\otimes k}\longrightarrow\mathbb Z,
\]
并定义\emph{由 $I$ 标号的\BookRecordIndex{T0066}标记点固定的不变量}为
\[
  \mathrm{GW}^{M,I}_{A,k}
  :=\widetilde{\mathrm{GW}}^{M,I}_{A,k}:
  H^*(M)^{\otimes k}\longrightarrow\mathbb Z.
\]
这里 $\widetilde{\mathrm{GW}}^{M,I}_{A,k}$ 如定理~\ref{thm:7-3-1} 所定义。
\end{definition}

\begin{remark}\label{rem:7-3-8}
当 $k<3$ 且 $A\ne0$ 时，不能直接从 $\widetilde{\mathrm{GW}}$ 恢复第~\ref{sec:7-1}节的
不变量 $\mathrm{GW}$；参见习题~\ref{ex:7-3-4}。另一方面，只要
$\mathrm{GW}^M_{A,k}$ 有定义，$k<3$ 且 $A\ne0$ 时就可借助除子公理从三点不变量恢复它
（见习题~\ref{ex:7-1-5}）。因此，对 $A\ne0$，可如下\BookRecordIndex{T0348}相容地定义零点、一点和两点不变量。
给定非零类 $A\in H_2(M;\mathbb Z)$，选取 $h\in H^2(M;\mathbb Z)$ 使
\[
  h(A)=\int_A h\ne0,
\]
并定义
\[
  \begin{aligned}
  \mathrm{GW}^M_{A,0}
  &:=\frac{1}{h(A)^3}
    \widetilde{\mathrm{GW}}^{M,\{1,2,3\}}_{A,3}(h,h,h),\\
  \mathrm{GW}^M_{A,1}(a)
  &:=\frac{1}{h(A)^2}
    \widetilde{\mathrm{GW}}^{M,\{1,2,3\}}_{A,3}(a,h,h),\\
  \mathrm{GW}^M_{A,2}(a,b)
  &:=\frac{1}{h(A)}
    \widetilde{\mathrm{GW}}^{M,\{1,2,3\}}_{A,3}(a,b,h).
  \end{aligned}
\]
特别地，当 $M$ 是满足 $c_1=0$ 的六维流形时，这给出了一种计数给定同调类中所有曲线
（包括\BookRecordIndex{T0113}多重覆盖曲线）数目的方法。（见第~11.3.4 节。）当 $A=0$ 时，$0\le k<3$ 的不变量
$\mathrm{GW}^M_{0,k}$ 与 $\mathrm{GW}^{M,I}_{0,k}$ 已在例~\ref{ex:7-1-3}、
例~\ref{ex:7-3-5} 和习题~\ref{ex:7-3-4} 中定义。它们全都为零，只有
$\#I=k\ne0$ 时例外；这时它们由杯积给出。
\end{remark}

\begin{remark}\label{rem:7-3-9}
在本书的解析框架内，不变量 $\mathrm{GW}^M$ 和 $\mathrm{GW}^{M,I}$ 也可以对某些非\BookRecordIndex{T0058}半正\BookRecordIndex{T0358}辛
流形 $(M,\omega)$ 定义。我们可以改为假设存在
$J_0\in\mathcal J_\tau(M,\omega)$，它在如下意义下是 $\kappa$-\BookRecordIndex{T0058}半正的：每个满足
$\int u^*\omega\le\kappa$ 的 $J_0$-全纯球面 $u:S^2\to M$ 都有非负 Chern 数；然后用
\BookRecordIndex{T0337}伪循环 $\operatorname{ev}_{J_0}$ 如前定义不变量。正如注~\ref{rem:6-6-2} 所指出的，所得
不变量此时可能依赖于\BookRecordIndex{T0088}殆复结构 $J_0$，因为被 $\omega$ \BookRecordIndex{T0372}驯服的 $\kappa$-\BookRecordIndex{T0058}半正\BookRecordIndex{T0088}殆复结构空间
可能不连通。因此，我们必须改变看法：把\BookRecordIndex{T0088}殆复结构 $J_0$ 而非\BookRecordIndex{T0363}辛形式 $\omega$ 视为给定的
数据。所以，这一推广在 凯勒 情形中特别有意义，因为此时首要数据是复结构，而非\BookRecordIndex{T0363}辛
形式。\footnote{与注~\ref{rem:7-1-11} 一样，这些顾虑事实上并无必要；\BookRecordIndex{T0370}虚模循环的存在表明，
这些不变量其实并不依赖于 $J_0$ 的选择。}下面精确陈述这一结果；它直接来自
注~\ref{rem:6-6-2}。

设 $(M,\omega)$ 是紧\BookRecordIndex{T0358}辛流形，且 $J_0\in\mathcal J_\tau(M,\omega)$ 关于 $\omega$ 是
$\kappa$-\BookRecordIndex{T0058}半正的。设 $A\in H_2(M;\mathbb Z)$ 满足 $\omega(A)\le\kappa$。则在 $M$ 上
所有\BookRecordIndex{T0088}殆复结构所成的空间中，存在 $J_0$ 的邻域 $\mathcal N(J_0)$，使同态
\[
  \mathrm{GW}^M_{A,k}=\mathrm{GW}^{M;J}_{A,k}:
  H^*(M^k)\longrightarrow\mathbb Z
\]
可由
\[
  \mathrm{GW}^{M;J}_{A,k}(a)
  =f\mathbin{\boldsymbol\cdot}\operatorname{ev}_J,
  \qquad
  \operatorname{ev}_J:\mathcal M^*_{0,k}(A;J)\longrightarrow M^k
\]
良定义，其中 $f:U\to M^k$ 是任意对偶于 $a$ 的\BookRecordIndex{T0337}伪循环；并且此同态与定义中所用的
$J\in\mathcal J_{\mathrm{reg}}(M,\omega)\cap\mathcal N(J_0)$ 的选择无关。只要
$\kappa\ge\omega(A)$，它也与\BookRecordIndex{T0372}驯服形式 $\omega$ 无关。它只依赖于 $J_0$ 在
$\kappa$-\BookRecordIndex{T0058}半正（殆）复结构族中的形变类。
\end{remark}

\begin{remark}\label{rem:7-3-10}
尽管我们几乎始终使用系数 $\mathbb Z$ 或 $\mathbb Q$，在应用中有时也适合忘掉定向而模
$2$ 计数曲线。在这种情形，与注~\ref{rem:6-5-8} 一样，应把记号 $H_*(M),H^*(M)$ 分别
理解为同调群 $H_*(M;\mathbb Z/2\mathbb Z)$ 与
$H^*(M;\mathbb Z/2\mathbb Z)$。前两节几乎所有结果都可显然地改写到这一设置中。例如，
定义~\ref{def:7-3-7} 的\BookRecordIndex{T0143}格罗莫夫--威滕不变量此时是同态
\[
  H^*(M;\mathbb Z/2\mathbb Z)^{\otimes k}
  \longrightarrow\mathbb Z/2\mathbb Z.
\]
不过，在某些情形中使用有理系数至关重要，例如像注~\ref{rem:7-3-8} 那样计数\BookRecordIndex{T0113}多重覆盖
曲线时。
\end{remark}
